\chapter{Evaluation and Results}
\label{ch:evaluation}
\vspace{2em}

% TODO: Introduce the chapter -- state what is being evaluated and the
% research questions it answers.

\section{Experimental Setup}
\label{eval:sec:setup}

% TODO: Dataset / course cohort, baselines, hardware, and evaluation
% protocol.

\section{Evaluation Metrics}
\label{eval:sec:metrics}

% TODO: Metrics used (e.g. answer accuracy, latency, faculty time saved,
% student satisfaction).

\subsection{Instructor-Hours Returned}
\label{eval:subsec:hours}

The headline workload metric promised in \autoref{review:subsec:gap} is
defined as follows. A student query is counted as \emph{deflected} when the
assistant answers it from the course corpus without instructor involvement
--- routine clarification, content-navigation, and logistics questions that
would otherwise reach the teaching staff. The estimated instructor-hours
returned per course per week is then

\begin{equation}
  \label{eval:eq:hours}
  H_{\mathrm{returned}}
    \;=\;
  \frac{1}{W} \sum_{q \,\in\, Q_{\mathrm{deflected}}} t_{\mathrm{handle}}(q),
\end{equation}

where $Q_{\mathrm{deflected}}$ is the set of deflected queries logged over an
observation window of $W$ weeks and $t_{\mathrm{handle}}(q)$ is the time an
instructor or teaching assistant would have spent answering query $q$
manually. Routine queries are assigned a deliberately conservative handling
time of 2--5 minutes, so $H_{\mathrm{returned}}$ is reported as a range whose
lower bound understates the true relief (it excludes context-switching costs
and out-of-hours response burden). Because every interaction is logged by the
deployed system, the metric reduces to a single aggregation over the
assistant's query logs.

% TODO: implement the aggregation as a PostgreSQL query over the production
% bot logs and report per-course results as a table in
% \autoref{eval:sec:results}.

A secondary, complementary signal is question-topic clustering over the same
logs: the distribution of what students actually ask provides an objective,
behaviour-based teaching signal that complements student evaluations of
teaching, whose limitations (expected-grade bias, perceived rather than
actual learning) were noted in \autoref{review:subsec:quality}.

% TODO: decide whether topic clustering is reported in this thesis or listed
% as future work in \autoref{ch:concl}.

\subsection{Token, Cost, Latency, and Answer-Quality Metrics}
\label{eval:subsec:tokenmetrics}

The compression stage is evaluated on the comparison dimensions the
prompt-compression literature itself uses
(\autoref{review:subsec:compression}), so that the evaluation speaks the same
metric language as the problem statement
(\autoref{intro:sec:problem}) and the baseline token audit
(\autoref{design:subsec:tokenaudit}):

\begin{enumerate}[nosep]
  \item \textbf{Input tokens per query}, before versus after compression,
        reported with the compression ratio --- the direct measure of the
        token inflation identified in \autoref{review:subsec:tokens};
  \item \textbf{Cost per query and per student per semester}, computed for
        the cloud-API analytics path; for the locally hosted generation path,
        where no per-token charge exists, latency and hardware-utilisation
        proxies are reported instead;
  \item \textbf{End-to-end latency} (retrieval $\rightarrow$ compression
        $\rightarrow$ generation), including the compression stage's own
        overhead --- compression adds an inference step, and its cost is
        reported rather than netted out, following the practice
        of~\citet{facl25:Jeong}~\cite{facl25:Jeong};
  \item \textbf{Answer quality} on course-material question answering, to
        verify that compression does not degrade answers --- mirroring how
        \textsc{Recomp} and TCRA-LLM report the compression--quality
        trade-off~\cite{iclr24:Xu, femnlp23:Liu}.
\end{enumerate}

The per-student-per-semester projection under dimension 2 --- baseline versus
compressed, extrapolated to a realistic course enrolment and set beside the
US\$7.50 per student per course per semester observed for
Ethel~\cite{tpt24:Kortemeyer} --- connects the evaluation directly back to
the institutional cost evidence of \autoref{intro:sec:problem}.

% TODO: report latency as P50/P95 once a larger probe set exists; the
% three-question probe in \autoref{eval:sec:results:compression} covers
% dimensions 1 and 3 only.
% TODO: decide the course-material QA evaluation set (n questions over the
% ingested Canvas corpus) for dimension 4.
% TODO: a first per-student-per-semester projection (baseline vs compressed) for
% the hosted generation path is in \autoref{eval:sec:results:providers}, built on
% ASSUMED usage rates. Replace them with rates measured from production logs, and
% chart the result against the Ethel benchmark.

\section{Results}
\label{eval:sec:results}

% TODO: Quantitative and qualitative results, with tables and figures.

\subsection{Context Compression: Token Reduction and Latency}
\label{eval:sec:results:compression}

The compression stage of \autoref{impl:sec:compression} was evaluated end-to-end
against the uncompressed RAG baseline on the deployed stack: identical questions,
identical retrieved chunks, generation by Llama~3.2 at temperature~0, with prompt
and output token counts taken directly from the inference runtime
(\texttt{prompt\_eval\_count}). The probe used the ingested CS3241 (computer
graphics) corpus with two in-syllabus questions --- Q1 on shadow rays, Q2 on
ray--sphere intersection --- and one deliberately out-of-syllabus question (Q3, on
TCP congestion control), with compression parameters $\tau = 0.5$, $k = 6$.
\autoref{eval:tab:compression} summarises the results.

\begin{table}[htb]
  \centering
  \caption{Uncompressed RAG versus extractive compression ($\tau=0.5$, $k=6$) on
           the deployed system. Compressed latency includes the compression
           overhead itself (1.1--2.0\,s per query on the 8\,GB host).}
  \label{eval:tab:compression}
  \begin{tabular}{l r r r}
    \hline
    \textbf{Metric} & \textbf{Raw RAG} & \textbf{Compressed} & \textbf{Change} \\
    \hline
    Prompt tokens, Q1            & 769            & 257           & $-66.6\%$ \\
    Prompt tokens, Q2            & 545            & 195           & $-64.2\%$ \\
    Response tokens, Q1          & 354            & 122           & $-65.5\%$ \\
    Response tokens, Q2          & 91             & 86            & $-5.5\%$  \\
    End-to-end latency, Q1       & 17.2\,s        & 7.7\,s        & $-55\%$   \\
    End-to-end latency, Q2       & 5.7\,s         & 5.2\,s        & $-9\%$    \\
    Context retained (chars), Q1 & 100\%          & 26.9\%        & ---       \\
    Context retained (chars), Q2 & 100\%          & 18.7\%        & ---       \\
    Off-syllabus context, Q3     & $\sim$446 tok. & 0 (empty)     & ---       \\
    \hline
  \end{tabular}
\end{table}

Three observations follow. First, compression removed roughly two-thirds of the
prompt tokens on in-syllabus questions while retaining the sentences that answer
them (the top-scoring sentence for Q1 was the corpus sentence defining shadow
rays, at cosine similarity 0.66). Because a compute-bound local model's latency
scales with prompt length, this translated into a net end-to-end speed-up
despite the added compression step. Second, the shorter, more focused context
also changed generation behaviour: on Q1 the baseline produced a longer answer
that drifted into adjacent material, whereas the compressed run stayed on the
retrieved definition, consistent with the distraction effect reported
by~\citet{tacl24:Liu}~\cite{tacl24:Liu}. Third, selective augmentation behaved as
designed: for the out-of-syllabus Q3, every candidate sentence scored below
$\tau$ (0.40--0.42, against 0.57--0.82 for in-syllabus sentences), so no context
was injected and the model correctly reported the topic as outside the course
materials; at a looser $\tau = 0.4$ the same query leaked three marginal
sentences, illustrating that $\tau$ is the precision--recall control of the
stage.

Compression is not uniformly beneficial, however. On Q2 the sentence containing
the ray--sphere intersection \emph{equation} itself scored below the retention
cut-off, and the compressed answer, while correct about the discriminant, could
no longer reproduce the equation --- an answer-quality regression directly
attributable to over-compression. This mirrors the recall risk inherent to
extractive compression~\cite{iclr24:Xu} and motivates the conservative default
($\tau = 0.5$ with $k = 6$) and the per-deployment tunability noted in
\autoref{impl:sec:compression}. These results are indicative rather than
statistically robust: they derive from a three-question probe over a single
course corpus, and a larger question set across multiple courses is required
before the reductions can be claimed as general.

\subsection{Local versus Hosted Inference}
\label{eval:sec:results:providers}

The deployment decision of \autoref{design:sec:deployment} was taken on
measurements of the two inference regimes on the same system, prompts and tasks.

\paragraph{Setup.} The agent loop was evaluated with a harness that replicates
the production loop: the same tool definitions, the same persona, a five-pass
cap, and tools withheld on the final pass. Tool results are stubbed and
byte-identical for every model, so differences reflect the model rather than live
data. The headline scenario is a multi-part request (\emph{``How am I doing on my
quizzes so far, and then quiz me on one weak topic?''}), run $n=10$ per model.

\begin{itemize}[nosep]
  \item \emph{Local models} (Llama~3.2 3B, Qwen3 14B, Q4 quantisation, thinking
        disabled) were measured on a 32\,GB Apple M1 Max through Ollama.
  \item \emph{Hosted models} (gpt-oss-20b and gpt-oss-120b, reasoning effort
        low) were measured on Groq's free tier from a client in Singapore.
\end{itemize}

One rubric scored both regimes. A turn counts as \emph{grounded} if the answer
states the stubbed performance score and names a stubbed weak topic. This
automatic scoring reproduces the original manual grounding scores for the local
models. Quiz generation was evaluated separately: the production
multiple-choice prompt, $n=10$ per model across five computing topics, with every
generated answer key checked by hand.

\begin{table}[htb]
  \centering
  \caption{Multi-part agent turn, local versus hosted inference ($n=10$, medians).
           Cost is marginal inference cost at list prices~\cite{groq26:Models}.}
  \label{eval:tab:providers}
  \begin{tabular}{l l r r r r}
    \hline
    \textbf{Model} & \textbf{Regime} & \textbf{Latency} & \textbf{Input tok.} & \textbf{Grounded} & \textbf{US\$/turn} \\
    \hline
    Llama~3.2 3B                  & local  & 3.3\,s  & 1{,}156 & 8/10  & ---     \\
    Qwen3 14B                     & local  & 11.2\,s & 1{,}595 & 10/10 & ---     \\
    gpt-oss-20b                   & hosted & 1.7\,s  & 2{,}248 & 10/10 & 0.00025 \\
    gpt-oss-120b\textsuperscript{a} & hosted & 2.2\,s  & 2{,}277 & 10/10 & 0.00049 \\
    \hline
  \end{tabular}\\[2pt]
  {\footnotesize \textsuperscript{a}With nullable optional tool parameters. With
  the original schema, all 10 turns were rejected by the provider's tool-argument
  validation. Input-token counts are not comparable across model families, because
  each family uses a different tokenizer.}
\end{table}

\begin{table}[htb]
  \centering
  \caption{Quiz question generation with the production prompt ($n=10$ per
           model). ``Valid'' counts outputs that parse to the required schema;
           ``Correct key'' counts questions, judged by hand, whose marked answer
           is right and unique.}
  \label{eval:tab:quiz-providers}
  \small
  \begin{tabular}{l l r r r r}
    \hline
    \textbf{Model} & \textbf{Regime} & \textbf{Latency} & \textbf{Valid} & \textbf{Correct key} & \textbf{US\$/question} \\
    \hline
    Llama~3.2 3B  & local  & 1.76\,s & 8/10  & 3/10 & ---      \\
    gpt-oss-20b   & hosted & 0.64\,s & 10/10 & 6/10 & 0.000074 \\
    gpt-oss-120b  & hosted & 0.86\,s & 10/10 & 8/10 & 0.000145 \\
    \hline
  \end{tabular}
\end{table}

\paragraph{Latency and quality.} \autoref{eval:tab:providers} and
\autoref{eval:tab:quiz-providers} give four results.
\begin{itemize}[nosep]
  \item \emph{Speed.} Hosted generation was 1.5--6.6 times faster on the agent
        turn. Generation throughput was roughly 475--965 tokens per second
        hosted, against 19--75 locally.
  \item \emph{Grounding.} Both hosted models matched the best local model,
        reporting the student's actual figures in 10 of 10 turns.
  \item \emph{Grounding versus tool use.} As in the earlier local benchmark,
        tool-invocation accuracy and grounding diverge. Qwen3 14B never called
        the question-generation tool, writing its own question instead, yet
        grounded every turn. Grounding fidelity, not tool accuracy, is therefore
        the discriminating metric for an assistant that reports grades.
  \item \emph{Quiz quality.} Schema validity overstated quiz quality for every
        model. Local Llama~3.2 produced parseable questions 8 times in 10, but
        correct answer keys only 3 times (for instance, it keyed ``a fixed rate
        regardless of congestion'' as TCP's congestion response). Hosted
        gpt-oss-120b produced correct keys 8 times in 10.
\end{itemize}

\paragraph{Provider-side validation.} The hosted provider validates tool-call
arguments against the declared JSON schema, which the local runtime never did.
This exposed a latent defect. When the student named no course, gpt-oss-120b
passed \texttt{null} for the optional course parameter, and the provider rejected
all 10 multi-part turns. Declaring optional parameters nullable reduced the
failures to 0 of 10 with 10 of 10 turns grounded, a result confirmed afterwards on
the production deployment. The local runtime had silently accepted the malformed
calls throughout development. A third hosted model, Qwen3.8 27B, was excluded: it
produced malformed tool calls in 4 of 13 tool-using turns and costs 5--13 times as
much per token.

\paragraph{Tokens in an agent loop.} The hosted results sharpen the thesis's
central claim that input tokens are the controllable cost driver. A tool-using
turn re-sends the system prompt, conversation, tool schemas and any retrieved
context on every pass, and the hosted models took three passes on the multi-part
turn. Retrieved context injected once per turn is therefore billed once per
\emph{pass}. For 1{,}000 tokens of retrieved context, that is 3{,}000 input
tokens. Adding that context raises the gpt-oss-20b turn from roughly US\$0.00025
to \$0.00049. Compressing it by 65\%, the reduction observed in
\autoref{eval:sec:results:compression}, brings it back to \$0.00034, a saving of
about 30\% per turn. Compression is thus worth proportionally more in agentic RAG
than in single-shot RAG. In the local regime, the same tokens appear instead as
prefill latency and, more acutely, as memory pressure on the host.

\paragraph{Cost per student.} \autoref{eval:tab:cost-projection} projects the
marginal generation cost per student per semester. It assumes 20 chat turns, 10
generated quiz questions and one feedback note per week over a 13-week semester.
These are round-number assumptions, not measured usage.

\begin{table}[htb]
  \centering
  \caption{Projected generation cost per student per semester (US\$), under
           assumed usage. ``Harness'' prompts carry no retrieved context; ``+RAG''
           adds about 1{,}000 tokens of retrieved context per pass; ``+RAG+comp.''
           compresses that context by 65\%. Model names abbreviate gpt-oss-20b
           and gpt-oss-120b.}
  \label{eval:tab:cost-projection}
  \begin{tabular}{l r r r}
    \hline
    \textbf{Model split} & \textbf{Harness} & \textbf{+RAG} & \textbf{+RAG+comp.} \\
    \hline
    Agent 20b, quiz 120b (deployed) & 0.090 & 0.148 & 0.110 \\
    Agent 120b, quiz 20b            & 0.141 & 0.258 & 0.182 \\
    \hline
  \end{tabular}
\end{table}

For the deployed split, the projection of US\$0.09--0.15 per student per semester
is roughly 50--85 times below the US\$7.50 per student per course per semester
reported for Ethel~\cite{tpt24:Kortemeyer}. The comparison is indicative only.
Ethel used a different class of model, different usage patterns and 2024 pricing.
Moreover, the free tier's limit of 8{,}000 tokens per minute per model allows only
about 1.4 retrieval-augmented agent turns per minute, so a production deployment
would need a paid tier whose pricing may differ.

\paragraph{Threats to validity.}
\begin{itemize}[nosep]
  \item \emph{Sample size.} $n=10$ per model supports ranking models, not
        fine-grained or tail-latency (P95) claims.
  \item \emph{Timing of the runs.} The local and hosted agent runs were taken
        five weeks apart on different runtime versions, with identical prompts
        and stubs.
  \item \emph{Client location.} Hosted latency was measured from Singapore and
        includes trans-Pacific round trips. Calls from the production server in
        the United States took 0.23--0.35\,s each, two to five times faster.
  \item \emph{Quiz rating.} Correctness was judged by a single, unblinded rater.
  \item \emph{Stubbed tools.} The harness exercises model behaviour rather than
        live retrieval and database access. An end-to-end timing of production
        turns on the hosted path remains to be done.
\end{itemize}

\section{Discussion}
\label{eval:sec:discussion}

% TODO: Interpret the results against the problem statement and threats to
% validity.

\section{Summary}
\label{eval:sec:summary}

% TODO: Recap the key findings.
